\chapter{Post-hoc Decisions and Protocol Deviations}\label{app:deviations}
Pre-registration is only as informative as its honesty about what was decided afterwards. This appendix lists, in time order, every decision that could have been influenced by data already seen, every addition to locked code, and every delegated choice. Entries refer to the repository's drift log (DR) and decision log (D).

\begin{enumerate}
\item \textbf{Crossing rule (D-018, DR-004; 22 Oct).} The target-reaching rule was changed from first crossing to two consecutive evaluations after the first tuning results had been seen, but before the lock of \code{prereg-v1} and before any final run. The reason (noise of single evaluations) is the same for every arm, and the re-scoring under three rules is recorded; a reader can still object.
\item \textbf{H6 generated from an exploratory lead (D-019, DR-005; 23 Oct).} Deep supervision was ahead of the baseline at the selection level; the hypothesis was added to the first pre-registration before any final seed, and its origin is stated.
\item \textbf{Status lines changed after the first results (DR-006).} The status lines of the locked hypothesis file were updated on 23 Oct, after the results of Study~1; these edits recorded outcomes and changed no rule, margin, threshold or seed set of any hypothesis, and further analyses were added after the results were seen.
\item \textbf{New arm added to locked code (DR-007; DR-010; DR-014; DR-015; DR-016).} The selective arm, the dose-response arm, Selective-Backprop and the options for the exploratory rounds were added to \code{src/} after the first lock, as additive changes with regression checks that the earlier arms' behaviour did not change. A first regression check of the current code against an earlier run used a different torch thread setting from the runs and was repeated with one thread; both attempts are logged (E-011).
\item \textbf{Shift rule of Studies 2 and 3 (D-023; 27 Oct).} The rule that the gain must exceed the re-tuning shift was changed from the absolute-change rule (written first) to \code{max\_gain}, by the assistant at the researcher's delegation, \emph{after} seeing that the second tuning A2 alone was $+34\%$ worse than A1 on the Study~1 seeds. The reason given is that a worse re-tuning does not weaken a claim of the form ``S beats the baseline''. The choice makes the shift smaller; the rule is stated in the lock and the other two rules are reported descriptively. On the actual seeds the verdicts would not have changed under the absolute-change rule: in Study~2 the gain of S30 ($35.1\%$ and $38.4\%$) exceeds the change of A2 ($26.2\%$ and $13.4\%$), and in Study~3 the gain of S20w0 ($48.5\%$ and $48.8\%$) exceeds that of A2 ($32.4\%$ and $11.3\%$).
\item \textbf{Interval level of Study 3 (D-026; 30 Oct).} The interval level $98.3\%$ ($=1-0.05/3$) for three comparisons was set by the assistant and recorded as an assistant decision. It is one constant in the code, and the verdicts do not change at $97.5\%$ (all intervals are far from their margins).
\item \textbf{Post-hoc contrast of S20 against S30 (E-023, F-020).} The contrast was computed on the seeds of the confirmatory test of Study~2 after the results were known, with the reading rule fixed before the run. It is exploratory and was the starting point of the Improve track.
\item \textbf{Choice of the primary arm of Study 3.} S20w0 was selected as the primary arm because it was the strongest lead of the exploratory rounds, among 14 arms; the confirmatory test used fresh seeds.
\item \textbf{H9 combines two changes.} The advantage over S30 is attributed jointly to the kept fraction and the warm-up; the split is exploratory.
\item \textbf{Selective-Backprop settings.} $\beta=2$, the history length and the one-epoch priming were fixed in advance by the researcher; the arm was tuned once.
\item \textbf{Documentation lag.} The status lines of H8--H10 in the repository were not yet updated at the time of writing (\cref{app:reproducibility}); the verdicts used here are recomputed from the raw logs with the locked analysis code.
\end{enumerate}
